% ============================================================================
% APPENDIX RB: LIT-BENABOU - Roland Bénabou Research Integration
% ============================================================================
% Category: LIT (Literature Integration)
% Version: 1.3
% Last Updated: 2026-02-04
% Dependencies: JT (LIT-TIROLE), K (LIT-FEHR), AU (CORE-AWARE)
% Language: english
% ============================================================================

% ----------------------------------------------------------------------------
% HEADER BLOCK
% ----------------------------------------------------------------------------
% Code: RB
% Category: LIT (Literature Integration)
% Title: Roland Bénabou Research Integration
% Author: EBF Framework
% Status: Complete
% Priority: High
% Evidence-Base: 30 papers (see data/paper-sources.yaml)
% ----------------------------------------------------------------------------

\chapter{LIT-BENABOU: Roland Bénabou Research Integration}
\label{app:lit-benabou}


\begin{tcolorbox}[colback=blue!5!white, colframe=blue!75!black,
    title=Quick Reference: Appendix RB]

\textbf{Appendix:} RB (LIT)

\textbf{Core Question:} What are the key findings in this domain?

\textbf{Key Concepts:}
\begin{itemize}[nosep]
\item Framework integration
\item Parameter validation
\item Cross-references
\end{itemize}

\textbf{Cross-References:} See related appendices below
\end{tcolorbox}

% ============================================================================
% CROSS-REFERENCE MAP
% ============================================================================
\begin{tcolorbox}[colback=blue!5!white,colframe=blue!75!black,title=Cross-Reference Map]
\textbf{Outgoing References:}
\begin{itemize}
    \item CORE-WHO (AAA): Multi-level welfare and social stratification
    \item CORE-WHAT (C): Identity, self-image, and belief utility
    \item CORE-HOW (B): Complementarity between intrinsic and extrinsic motivation
    \item CORE-WHEN (V): Context effects on motivation and beliefs
    \item CORE-WHERE (BBB): Parameters for motivation crowding
    \item CORE-AWARE (AU): Awareness, self-deception, and motivated reasoning
    \item CORE-READY (AV): Willingness, self-confidence, and action thresholds
    \item CORE-STAGE (AW): Behavioral change and identity transformation
    \item LIT-TIROLE (JT): Behavioral IO collaboration
    \item LIT-FALK (FAL): Moral reasoning and prosocial behavior collaboration
    \item LIT-FEHR (K): Social preferences comparison
    \item LIT-KAHNEMAN (L): Dual-system foundations
    \item DOMAIN-EDU (AB): Education and human capital
    \item DOMAIN-POL (AG): Political economy of redistribution
\end{itemize}

\textbf{Incoming References:}
\begin{itemize}
    \item LIT-TIROLE (JT): Joint behavioral IO research
    \item DOMAIN-ORG (AJ): Organizational behavior and groupthink
    \item PREDICT-MOTIVATION (AS): Motivation crowding predictions
\end{itemize}
\end{tcolorbox}

% ============================================================================
% CHAPTER LINKAGE
% ============================================================================
\begin{tcolorbox}[colback=green!5!white,colframe=green!75!black,title=Chapter Linkage]
\textbf{Primary Chapter:} Chapter 11 (Beliefs and Identity)

\textbf{Supporting Chapters:}
\begin{itemize}
    \item Chapter 5: Self-control and time preferences
    \item Chapter 7: Social stratification and inequality
    \item Chapter 10: Motivation and incentive design
    \item Chapter 14: Education policy applications
\end{itemize}
\end{tcolorbox}

% ============================================================================
% ABSTRACT AND QUICK REFERENCE
% ============================================================================
\begin{abstract}
This appendix integrates Roland Bénabou's foundational contributions to behavioral economics and political economy into the EBF framework. Bénabou's research program spans three interconnected domains: (1) the psychology of motivation, self-confidence, and identity (primarily with Tirole), (2) inequality, social mobility, and redistribution preferences, and (3) collective beliefs, groupthink, and organizational delusions. His work provides essential theoretical foundations for understanding how beliefs serve as assets, why incentives can backfire, and how inequality perpetuates through educational stratification. The EBF framework incorporates Bénabou's insights through seven axioms (RB-1 to RB-7) that formalize motivated beliefs, self-confidence dynamics, the POUM hypothesis (Prospect of Upward Mobility), groupthink mechanisms, moral narratives and imperatives, and the local complementarities that drive educational segregation.
\end{abstract}

\begin{tcolorbox}[colback=gray!5!white,colframe=gray!75!black,title=Quick Reference]
\textbf{Core Contributions:}
\begin{itemize}
    \item Motivation psychology: Intrinsic/extrinsic interaction, self-confidence, willpower
    \item Identity economics: Beliefs as assets, moral reasoning, taboos
    \item Inequality and growth: Education, stratification, local complementarities
    \item Political economy: POUM hypothesis, redistribution preferences
    \item Organizational behavior: Groupthink, collective delusions, denial
    \item Religion and innovation: Science-religion dynamics
\end{itemize}

\textbf{Key Parameters:}
\begin{itemize}
    \item Self-confidence multiplier: $\partial e / \partial \hat{\theta} > 0$ (effort increases with confidence)
    \item Motivation crowding threshold: Signal content of incentive determines direction
    \item POUM effect: $\mathbb{E}[y'] > y \Rightarrow$ oppose redistribution even if $y < \bar{y}$
    \item Groupthink cascade: $\Pr(\text{speak up}) \downarrow$ as perceived consensus $\uparrow$
\end{itemize}

\textbf{Evidence Tier:} Mixed (Tier 2-5)
\begin{itemize}
    \item Motivation crowding: Tier 1-2 (lab and field experiments)
    \item Inequality and growth: Tier 2-3 (cross-country, panel data)
    \item POUM hypothesis: Tier 2-3 (survey evidence, natural experiments)
    \item Groupthink: Tier 3-4 (case studies, organizational research)
\end{itemize}
\end{tcolorbox}

% ============================================================================
% FUNDAMENTAL QUESTION
% ============================================================================
% -----------------------------------------------------------------------------
% APPENDIX SCOPE BOX
% -----------------------------------------------------------------------------

\begin{tcolorbox}[
    colback=orange!5!white,
    colframe=orange!75!black,
    title={\textbf{Appendix Scope: Ziel / In-Scope / Out-of-Scope / Constraints}},
    fonttitle=\bfseries
]

\textbf{Ziel (Objective):}
\begin{itemize}[nosep]
    \item How do beliefs, identity, and social context shape economic behavior---and how do these psychological factors interact with inequality to perpetuate or transform social structures?
\end{itemize}

\textbf{In-Scope:}
\begin{itemize}[nosep]
    \item Fundamental Question
    \item Key Results and Evidence
    \item EBF Integration
    \item Worked Example: Volunteer Motivation
\end{itemize}

\textbf{Out-of-Scope (delegiert an andere Appendices):}
\begin{itemize}[nosep]
    \item REF-GLOSSARY $\rightarrow$ Appendix G
    \item REF-GLOSSARY $\rightarrow$ Appendix G
    \item LIT-TIROLE $\rightarrow$ Appendix JT
\end{itemize}

\textbf{Constraints (Anwendungsgrenzen):}
\begin{itemize}[nosep]
    \item [TODO: Define when this appendix does NOT apply]
    \item [TODO: Define required prerequisites]
    \item [TODO: Define known limitations]
\end{itemize}

\textbf{Lieferobjekte (Deliverables):}
\begin{itemize}[nosep]
    \item 4 Table(s)
    \item 12 Equation(s)
    \item 10 Axiom(s) (RB-1 to RB-10)
    \item Worked Example(s)
\end{itemize}

\end{tcolorbox}


\section{Fundamental Question}
\label{sec:rb-fundamental}

\begin{tcolorbox}[colback=red!5!white,colframe=red!75!black,title=Fundamental Question]
\textbf{How do beliefs, identity, and social context shape economic behavior---and how do these psychological factors interact with inequality to perpetuate or transform social structures?}
\end{tcolorbox}

Roland Bénabou's research addresses a fundamental gap between standard economics and observed behavior: people hold beliefs that seem ``wrong'' by Bayesian standards, they respond to incentives in unexpected ways, and they support policies that appear contrary to their material interests. His work shows that these apparent anomalies are rational responses when we account for:

\begin{enumerate}
    \item \textbf{Beliefs as consumption goods:} Beliefs provide direct utility (anticipatory pleasure, self-image maintenance)
    \item \textbf{Beliefs as investment goods:} Self-confidence affects motivation and performance
    \item \textbf{Social interdependencies:} Individual outcomes depend on community composition
    \item \textbf{Collective belief dynamics:} Organizations can develop shared delusions
\end{enumerate}

For the EBF framework, Bénabou's contributions are essential because they:
\begin{itemize}
    \item Formalize how awareness (CORE-AWARE) is endogenously shaped by motivated reasoning
    \item Explain when incentives enhance versus undermine intrinsic motivation
    \item Connect individual psychology to macro-level inequality dynamics
    \item Provide tools for understanding collective failures in organizations and societies
\end{itemize}

% ============================================================================
% THEORETICAL FOUNDATIONS
% ============================================================================
\section{Theoretical Foundations}
\label{sec:rb-theory}

\subsection{Self-Confidence and Personal Motivation}
\label{subsec:rb-confidence}

Bénabou and Tirole's self-confidence model \citep{benabou_tirole_2002_self_confidence} shows why people engage in self-deception:

\begin{equation}
V(\hat{\theta}) = \max_e \left[ \mathbb{E}_{\hat{\theta}}[u(e, \theta)] + \beta \cdot V(\hat{\theta}') \right]
\label{eq:rb-confidence}
\end{equation}

where $\hat{\theta}$ is self-assessed ability, $e$ is effort, and future self-confidence $\hat{\theta}'$ depends on current beliefs and outcomes.

\begin{axiom}[Self-Confidence as Motivation]
\label{ax:rb-1}
\textbf{RB-1:} Self-confidence has instrumental value---higher $\hat{\theta}$ increases effort and performance:
\begin{equation}
\frac{\partial e^*}{\partial \hat{\theta}} > 0, \quad \frac{\partial \mathbb{E}[y]}{\partial \hat{\theta}} > 0
\end{equation}
This creates demand for positive self-views. Rational agents may engage in selective memory, biased attribution, or avoidance of diagnostic information to maintain beneficial beliefs.
\end{axiom}

\textbf{EBF Integration:} Axiom RB-1 connects directly to CORE-AWARE (AU)---awareness is not just about information but about \textit{willingness} to process information that threatens self-image. It also connects to CORE-READY (AV)---self-confidence affects the perceived threshold for action.

\textbf{Key Insight:} There is an optimal level of self-deception. Too little confidence undermines motivation; too much leads to poor decisions from overconfidence.

\subsection{Intrinsic and Extrinsic Motivation}
\label{subsec:rb-motivation}

The Bénabou-Tirole motivation model \citep{benabou_tirole_2003_intrinsic} formalizes when incentives backfire:

\begin{equation}
e^* = \arg\max_e \left[ \underbrace{v \cdot a(e)}_{\text{intrinsic}} + \underbrace{\mu \cdot y(e)}_{\text{image}} + \underbrace{b \cdot e}_{\text{extrinsic}} - C(e) \right]
\label{eq:rb-motivation}
\end{equation}

\begin{axiom}[Motivation Crowding via Signaling]
\label{ax:rb-2}
\textbf{RB-2:} Extrinsic incentives affect effort through two channels:
\begin{equation}
\frac{de^*}{db} = \underbrace{\frac{\partial e^*}{\partial b}\bigg|_{\hat{v}}}_{\text{direct (+)}} + \underbrace{\frac{\partial e^*}{\partial \hat{v}} \cdot \frac{\partial \hat{v}}{\partial b}}_{\text{signal (+ or -)}}
\end{equation}
Crowding out occurs when incentives signal that the task is unpleasant ($\partial \hat{v}/\partial b < 0$) or that the principal distrusts the agent. Crowding in occurs when incentives signal task importance or confidence in the agent.
\end{axiom}

\textbf{EBF Integration:} Axiom RB-2 is central to EBF's intervention design. It connects to CORE-HOW (B) through the complementarity structure---intrinsic and extrinsic motivation can be complements or substitutes depending on what the incentive signals.

\subsection{Identity, Morals, and Taboos}
\label{subsec:rb-identity}

Bénabou and Tirole's identity model \citep{benabou_tirole_2011_identity} treats moral beliefs as strategic assets:

\begin{equation}
U = u(\text{actions}, \text{outcomes}) + \lambda \cdot \text{self-image}(\text{beliefs}, \text{actions})
\label{eq:rb-identity}
\end{equation}

\begin{axiom}[Beliefs as Assets]
\label{ax:rb-3}
\textbf{RB-3:} Agents invest in beliefs that:
\begin{enumerate}
    \item Provide anticipatory utility (pleasant to hold)
    \item Support a positive self-image
    \item Enable commitment to valued behaviors
    \item Create social bonding through shared identity
\end{enumerate}
Taboos emerge endogenously: explicitly considering trade-offs (e.g., pricing lives) destroys the ``sacred'' value that enables commitment.
\end{axiom}

\textbf{EBF Integration:} Axiom RB-3 explains why information campaigns often fail to change behavior---new information may threaten identity investments. It connects to CORE-STAGE (AW)---identity transformation is a critical part of the behavioral change journey.

\subsection{Inequality, Stratification, and Growth}
\label{subsec:rb-inequality}

Bénabou's inequality research \citep{benabou_1996_inequality_growth, benabou_1996_heterogeneity} shows how local complementarities create persistent stratification:

\begin{equation}
h_{i,t+1} = f(h_{i,t}, \bar{h}_{N(i),t}, e_{i,t})
\label{eq:rb-inequality}
\end{equation}

where $h$ is human capital, $\bar{h}_{N(i)}$ is average human capital in agent $i$'s neighborhood, and $e$ is educational investment.

\begin{axiom}[Local Complementarities and Stratification]
\label{ax:rb-4}
\textbf{RB-4:} When human capital production exhibits local complementarities ($\partial^2 f / \partial h_i \partial \bar{h}_N > 0$), stratified equilibria emerge:
\begin{enumerate}
    \item High-skill agents sort into high-skill communities
    \item This sorting increases inequality in the next generation
    \item Multiple steady states exist with different inequality levels
    \item Initial conditions (history) determine which equilibrium obtains
\end{enumerate}
Integrated communities may be more efficient but are unstable without policy intervention.
\end{axiom}

\textbf{EBF Integration:} Axiom RB-4 connects to CORE-WHO (AAA) through the welfare hierarchy---community-level welfare affects individual outcomes. It explains why individual-level interventions may fail when systemic stratification persists.

\subsection{The POUM Hypothesis}
\label{subsec:rb-poum}

Bénabou and Ok's POUM model \citep{PAP-benabou_ok_2001_mobility} explains why poor voters may oppose redistribution:

\begin{equation}
\text{Support redistribution} \Leftrightarrow \mathbb{E}[u(y'(1-\tau) + \tau\bar{y}')] > u(y(1-\tau) + \tau\bar{y})
\label{eq:rb-poum}
\end{equation}

\begin{axiom}[Prospect of Upward Mobility]
\label{ax:rb-5}
\textbf{RB-5:} Under concave utility and sufficient income mobility, agents below the mean may rationally oppose redistribution if they expect their future income to exceed the (future) mean:
\begin{equation}
y_t < \bar{y}_t \text{ but } \mathbb{E}[y_{t+1}] > \mathbb{E}[\bar{y}_{t+1}] \Rightarrow \text{oppose redistribution}
\end{equation}
This explains the ``American paradox''---high tolerance for inequality despite relatively low mobility---if beliefs about mobility are inflated.
\end{axiom}

\textbf{EBF Integration:} Axiom RB-5 connects motivated beliefs (RB-3) to political preferences. Beliefs about social mobility may be systematically biased upward to maintain self-image and hope.

\subsection{Groupthink and Collective Delusions}
\label{subsec:rb-groupthink}

Bénabou's groupthink model \citep{PAP-benabou_2013_groupthink} explains organizational failures:

\begin{equation}
\Pr(\text{speak up}) = g\left(\text{private signal}, \text{perceived consensus}, \text{reputational cost}\right)
\label{eq:rb-groupthink}
\end{equation}

\begin{axiom}[Collective Denial Cascades]
\label{ax:rb-6}
\textbf{RB-6:} Organizations can develop ``mutually assured delusion'' when:
\begin{enumerate}
    \item Individuals have reputational incentives to conform
    \item Early silence is interpreted as agreement
    \item Sunk costs create desire to believe past decisions were correct
    \item Dissent becomes increasingly costly as consensus appears stronger
\end{enumerate}
This explains corporate scandals, policy failures, and market bubbles where ``everyone knew'' but no one spoke up.
\end{axiom}

\textbf{EBF Integration:} Axiom RB-6 extends individual motivated beliefs (RB-3) to organizations. It connects to CORE-WHO (AAA)---organizational-level pathologies emerge from individual incentives.

\subsection{Moral Narratives and Imperatives}
\label{subsec:rb-narratives}

Bénabou, Falk, and Tirole's narratives model \citep{benabou_2018_narratives} analyzes how moral justifications and rules spread through society:

\begin{equation}
U = (ve - c/\beta)a + \mu \hat{v}(a)
\label{eq:rb-narratives}
\end{equation}

where $a \in \{0,1\}$ is the moral action, $v$ is true value, $e$ is externality, $c$ is cost, $\beta \leq 1$ captures self-control, and $\mu \hat{v}(a)$ is image utility based on inferred type.

\begin{axiom}[Moral Narratives and Strategic Transmission]
\label{ax:rb-7}
\textbf{RB-7:} Moral justifications (narratives) and categorical rules (imperatives) serve distinct functions:

\textbf{Narratives} are belief manipulations that affect the perceived parameters of moral decisions:
\begin{itemize}
    \item \textit{Positive narratives} (``responsibility'') emphasize $v$ or $e$, promoting prosocial behavior
    \item \textit{Negative narratives} (``excuses'') emphasize costs or uncertainty, enabling selfish behavior
\end{itemize}

\textbf{Key result on strategic complementarity:}
\begin{equation}
\frac{\partial^2 \Pi}{\partial n_i \partial n_j} \begin{cases} < 0 & \text{for negative narratives (substitutes)} \\ > 0 & \text{for positive narratives (complements)} \end{cases}
\end{equation}

\textbf{Imperatives} are categorical rules that bypass cost-benefit analysis:
\begin{itemize}
    \item Effective when $\mu$ (image concern) is high
    \item Require perceived moral authority of the source
    \item Enable commitment by removing deliberation
\end{itemize}
\end{axiom}

\textbf{Network Transmission:} The spread of narratives depends on network structure $\rho$ (persistence/homophily):
\begin{equation}
N^-_A = \frac{\pi(1-x)}{1 - \rho(1-x)}, \quad N^-_B = \frac{(1-\pi)x}{1 - \rho x}
\label{eq:rb-network}
\end{equation}

Key insight: More network mixing (lower $\rho$) increases exposure to prosocial agents, raising moral behavior overall.

\textbf{EBF Integration:} Axiom RB-7 connects to:
\begin{itemize}
    \item CORE-AWARE (AU): Narratives manipulate awareness of costs and benefits
    \item CORE-HOW (B): Narrative transmission exhibits strategic complementarity
    \item CORE-WHEN (V): Network structure $\rho$ is a key context moderator
    \item LIT-FALK (FAL): Joint work with Armin Falk on prosocial behavior
\end{itemize}

\textbf{Key Parameters from Paper:}
\begin{itemize}
    \item $\mu$ (image concern): Weight on reputation/self-image in utility
    \item $\beta$ (self-control): Present bias parameter, $\beta \leq 1$
    \item $\rho$ (network persistence): Correlation of types on network, $\rho \in [0,1]$
\end{itemize}

See Parameter Registry entries PAR-BEH-014 ($\mu$) and PAR-BEH-015 ($\rho$) for empirical calibrations.

\subsection{Ends versus Means: Context-Dependent Moral Reasoning}
\label{subsec:rb-ends-means}

Bénabou, Falk, and Henkel \citep{benabou_2024_ends_means} challenge the assumption of stable moral ``types'' (Kantian vs. Utilitarian). Using four novel \textbf{Ends-versus-Means (EVM)} games, they show:

\begin{axiom}[Context-Dependent Moral Reasoning]
\label{ax:rb-11}
\textbf{RB-11:} Moral reasoning weight $\theta$ is context-dependent, not a stable trait:
\begin{equation}
U_{\text{moral}} = \theta(\Psi) \cdot U_{\text{outcome}} + (1-\theta(\Psi)) \cdot U_{\text{rule}}
\label{eq:rb-moral-utility}
\end{equation}

\textbf{Critical Negative Result:}
\begin{equation}
\gamma(\text{EVM}_i, \text{EVM}_j) \approx 0 \quad \text{for } i \neq j
\end{equation}
Cross-context correlations are effectively zero (mean $r = 0.035$), contradicting stable type theories.

\textbf{Key Findings:}
\begin{itemize}
    \item \textbf{No stable types}: Same individual may be deontological in lying context (44\%) but consequentialist in group donation (80\%)
    \item \textbf{EVM $\perp$ Prosociality}: Moral reasoning is orthogonal to altruism ($r \approx 0.025$)
    \item \textbf{Rule-following $\neq$ Deontology}: Kimbrough-Vostroknutov $\kappa$ does not predict EVM choices
\end{itemize}
\end{axiom}

\textbf{Deontological Choice Rates by Context:}
\begin{center}
\begin{tabular}{lcc}
\toprule
\textbf{Context} & $\theta_{\text{deont}}$ & 95\% CI \\
\midrule
Lying Game & 0.44 & [0.39, 0.49] \\
Bribe Game & 0.38 & [0.33, 0.43] \\
Statement Choice & 0.28 & [0.23, 0.33] \\
Group Donation & 0.20 & [0.16, 0.24] \\
\bottomrule
\end{tabular}
\end{center}

\textbf{EBF Integration:} RB-11 provides a \textbf{critical counterexample} to the general pattern of cross-context complementarity:
\begin{itemize}
    \item \textbf{CORE-HOW (B)}: $\gamma \approx 0$ in moral domain contradicts typical $\gamma > 0$ patterns (see Chapter 5 Counterexample Box)
    \item \textbf{CORE-WHEN (V)}: $\theta = f(\Psi)$ demonstrates Parameter Context Transformation (PCT) in moral reasoning
    \item \textbf{CORE-AWARE (AU)}: $A_{\text{moral}}$ varies 20--44\% across contexts
    \item \textbf{CORE-WHO (AAA)}: Heterogeneity is context-driven, not person-driven
\end{itemize}

\textbf{Key Parameters:} See PAR-MR-001 to PAR-MR-004 for validated parameter values.

\textbf{Intervention Implications:}
\begin{itemize}
    \item Cannot assume stable Kantian vs. Utilitarian types when designing interventions
    \item Prosocial interventions $\neq$ moral reasoning interventions (orthogonal)
    \item Moral messaging must be context-specific, not person-specific
\end{itemize}

\subsection{Identity as Self-Image: A Unified Framework}
\label{subsec:rb-identity-self-image}

Bénabou and Henkel's comprehensive handbook chapter \citep{benabou_henkel_2025_identity_self_image} provides the \textbf{unified framework} for belief-based identity economics, synthesizing 20+ years of research:

\begin{equation}
U = u(a, \theta) + \mu \cdot v(\sigma) \quad \text{where} \quad \sigma = E[\theta|I]
\label{eq:rb-self-image}
\end{equation}

The framework identifies \textbf{three approaches} differing in timing of self-image utility:
\begin{enumerate}
    \item \textbf{Anticipatory}: Utility from anticipating future consumption (Caplin-Leahy 2001)
    \item \textbf{Contemporaneous}: Ego utility from current self-esteem
    \item \textbf{Retrospective}: Caring about future views of past selves
\end{enumerate}

\begin{axiom}[Self-Image Utility]
\label{ax:rb-8}
\textbf{RB-8:} Identity utility depends on self-perceived type:
\begin{equation}
U_{\text{IDN}} = \mu \cdot v(\sigma) \quad \text{where} \quad \sigma = E[\theta|I]
\end{equation}

\textbf{Building blocks:}
\begin{itemize}
    \item \textbf{Types} ($\theta$): Underlying characteristic (ability, morality, health)
    \item \textbf{Self-image} ($\sigma$): Agent's belief about own type
    \item \textbf{Image utility} ($v(\sigma)$): Intrinsic value of self-perceived type
    \item \textbf{Image concern} ($\mu$): Weight on identity in overall welfare
\end{itemize}
\end{axiom}

\begin{axiom}[Information Avoidance]
\label{ax:rb-9}
\textbf{RB-9:} Agents may strategically avoid information when ego utility exceeds instrumental value:
\begin{equation}
V(\text{info}) = V_{\text{instrumental}} - V_{\text{ego loss}}
\end{equation}

Agents avoid information when ego costs dominate, explaining:
\begin{itemize}
    \item Reluctance to get genetic/health tests (Oster et al. 2013)
    \item Not checking portfolios during downturns (Karlsson et al. 2009)
    \item Avoiding performance feedback (Kőszegi 2003)
\end{itemize}
\end{axiom}

\begin{axiom}[Fear of Asking]
\label{ax:rb-10}
\textbf{RB-10:} Help-seeking is deterred by fear of rejection when rejection reveals low valuation:
\begin{equation}
U_R(h, F) = wh + \psi(F) \quad \text{where} \quad \psi(\text{reject}) = -\lambda\alpha[\tilde{g} - \bar{g}_N] < 0
\end{equation}

The \textbf{asking threshold} $w^*$ is determined by rejection sensitivity:
\begin{equation}
w^* = \lambda\alpha[M^+_N(\hat{g}(w^*)) - \bar{g}_N]
\end{equation}

\textbf{Key Results} from Bénabou, Jaroszewicz \& Loewenstein (2022) \citep{benabou_2022_hurts_ask}:
\begin{itemize}
    \item \textbf{Non-monotonic asking}: Higher need $w$ can reduce asking when rejection costs dominate
    \item \textbf{Discouragement effect}: Non-offers discourage asking ($w^*_N > w^*_O$)
    \item \textbf{Waiting trap}: Pareto-inferior equilibrium where both parties wait
    \item \textbf{Offering vs asking}: When $\lambda\alpha$ high, offering norms dominate asking norms
\end{itemize}

\textbf{Applications:}
\begin{itemize}
    \item Welfare non-take-up despite eligibility (Currie 2006)
    \item Reluctance to ask colleagues for help (Flynn \& Lake 2008)
    \item Healthcare avoidance and diagnosis aversion
    \item Charitable giving solicitation effects
\end{itemize}
\end{axiom}

\textbf{Key Mechanisms:}

\begin{enumerate}
    \item \textbf{Self-Signaling} (Theory MS-IB-014): Actions reveal information about type to oneself:
    \begin{equation}
    \sigma' = E[\theta|a, I] \quad \text{(Bayesian updating on actions)}
    \end{equation}
    Creates \textit{diagnostic utility}---value of learning about oneself through choices.

    \item \textbf{Motivated Beliefs} (Brunnermeier-Parker 2005): Optimal belief distortion:
    \begin{equation}
    \max_{\pi} \left[ U(a^*(\pi), \pi) - \kappa \cdot D(\pi||\pi_0) \right]
    \end{equation}

    \item \textbf{Selective Memory}: Action-dependent recall $\rho(a) = P(\text{recall}|a)$

    \item \textbf{Dual-Self Models}: Self-control as game between selves with different $\beta$
\end{enumerate}

\textbf{EBF Integration:} The unified framework from RB-8, RB-9, and RB-10 connects to:
\begin{itemize}
    \item \textbf{CORE-WHO (AAA)}: Self-image $\sigma$ revealed by others' valuation; rejection signals low $g$
    \item \textbf{CORE-AWARE (AU)}: Information avoidance extended to asking decisions (fear of rejection)
    \item \textbf{CORE-WHAT (C)}: $\mu v(\sigma)$ and $\psi(F)$ as identity/image utility components
    \item \textbf{CORE-HOW (B)}: Asking-offering complementarity $\gamma(\text{offer}, \text{ask})$ depends on rejection costs
    \item \textbf{CORE-READY (AV)}: Asking threshold $w^*$ determines readiness to seek help
    \item Theory Catalog: MS-IB-014 (Self-Signaling), MS-IB-015 (Diagnostic Utility), MS-IB-016 (Information Avoidance), MS-IB-017 (Unified Framework), MS-IB-018 (Fear of Asking)
\end{itemize}

\textbf{Key Parameters:}
\begin{itemize}
    \item $\mu$: Image concern intensity [0, 1]
    \item $\sigma$: Self-perceived type (continuous)
    \item $\lambda_A$: Anticipatory utility weight [0, 1]
    \item $\rho$: Memory persistence [0, 1]
    \item $\kappa$: Self-deception cost [0, $\infty$)
    \item $\lambda_R$: Rejection sensitivity (RB-10) [1.5, 3.0]
    \item $\alpha$: Rejection scaling factor (RB-10) [0, 1]
    \item $w^*$: Asking threshold (RB-10) (endogenous)
    \item $\tilde{g}$: Offering threshold (RB-10) (endogenous)
\end{itemize}

% ============================================================================
% KEY RESULTS AND EVIDENCE
% ============================================================================
\section{Key Results and Evidence}
\label{sec:rb-results}

\subsection{Motivation Crowding Evidence}

\begin{table}[htbp]
\centering
\caption{Motivation Crowding: Empirical Evidence}
\label{tab:rb-crowding}
\begin{tabular}{llcc}
\toprule
\textbf{Study} & \textbf{Finding} & \textbf{RB-2 Support} & \textbf{Tier} \\
\midrule
Gneezy \& Rustichini (2000) & Small fines increased tardiness & Yes (signal) & 1 \\
Frey \& Jegen (2001) & Meta-analysis of 128 studies & Yes (conditional) & 2 \\
Ariely et al. (2009) & High incentives hurt performance & Yes (pressure) & 1 \\
Mellström \& Johannesson (2008) & Payment reduced blood donation & Yes (image) & 1 \\
Heyman \& Ariely (2004) & Money vs. gifts affect effort & Yes (framing) & 1 \\
\bottomrule
\end{tabular}
\end{table}

\subsection{Inequality and Stratification Evidence}

\begin{table}[htbp]
\centering
\caption{Local Complementarities and Stratification}
\label{tab:rb-stratification}
\begin{tabular}{llcc}
\toprule
\textbf{Study} & \textbf{Finding} & \textbf{RB-4 Support} & \textbf{Tier} \\
\midrule
Chetty et al. (2014) & Neighborhood effects on mobility & Yes & 1 \\
Card \& Rothstein (2007) & School segregation effects & Yes & 1 \\
Durlauf (2004) & Neighborhood sorting review & Yes & 2 \\
Aaronson (1998) & Housing segregation persistence & Yes & 2 \\
\bottomrule
\end{tabular}
\end{table}

\subsection{POUM and Redistribution Preferences}

\begin{table}[htbp]
\centering
\caption{Mobility Beliefs and Redistribution Preferences}
\label{tab:rb-poum}
\begin{tabular}{llcc}
\toprule
\textbf{Study} & \textbf{Finding} & \textbf{RB-5 Support} & \textbf{Tier} \\
\midrule
Alesina \& La Ferrara (2005) & Mobility beliefs predict preferences & Yes & 2 \\
Piketty (1995) & Dynastic learning model & Partial & 3 \\
Alesina \& Glaeser (2004) & US-Europe comparison & Yes & 3 \\
Chetty et al. (2017) & Biased mobility perceptions & Yes & 1 \\
\bottomrule
\end{tabular}
\end{table}

% ============================================================================
% EBF INTEGRATION
% ============================================================================
\section{EBF Integration}
\label{sec:rb-integration}

\subsection{CORE Framework Mapping}

\begin{table}[htbp]
\centering
\caption{Bénabou's Contributions Mapped to 10C Framework}
\label{tab:rb-9c-mapping}
\begin{tabular}{lll}
\toprule
\textbf{CORE} & \textbf{Bénabou Contribution} & \textbf{Application} \\
\midrule
WHO (AAA) & Local complementarities & Community-level welfare \\
WHAT (C) & Belief utility, self-image & Non-material utility components \\
HOW (B) & Motivation complementarities & Incentive design \\
WHEN (V) & Context-dependent signaling & When incentives backfire \\
WHERE (BBB) & Crowding parameters & Incentive effect estimation \\
AWARE (AU) & Motivated reasoning & Endogenous awareness \\
READY (AV) & Self-confidence effects & Action threshold dynamics \\
STAGE (AW) & Identity transformation & Behavioral change journey \\
\bottomrule
\end{tabular}
\end{table}

\subsection{Implications for Intervention Design}

\begin{tcolorbox}[colback=green!5!white,colframe=green!75!black,title=When Incentives Work]
\begin{itemize}
    \item Task is already perceived as instrumental (no intrinsic value to protect)
    \item Incentive signals confidence or importance
    \item Agent has little prior intrinsic motivation
    \item Context frames incentive as recognition, not control
\end{itemize}
\end{tcolorbox}

\begin{tcolorbox}[colback=red!5!white,colframe=red!75!black,title=When Incentives Backfire]
\begin{itemize}
    \item Task has strong intrinsic or identity value
    \item Incentive signals distrust or unpleasantness
    \item Small incentive implies task is not valuable
    \item Context frames incentive as controlling
\end{itemize}
\end{tcolorbox}

% ============================================================================
% WORKED EXAMPLE
% ============================================================================
\section{Worked Example: Volunteer Motivation}
\label{sec:rb-example}

\begin{tcolorbox}[colback=yellow!5!white,colframe=yellow!75!black,title=Worked Example: Blood Donation Campaign]
\textbf{Scenario:} A blood bank considers offering payment to increase donations. Current donation rate: 10\%.

\textbf{Standard Economic Prediction:}
Payment $b > 0$ should increase donations: $\partial D / \partial b > 0$

\textbf{Bénabou-Tirole Analysis (RB-2):}
\begin{equation}
D(b) = D_0 + \alpha \cdot b - \beta \cdot \mathbf{1}_{b > 0} \cdot \text{signal}(b)
\end{equation}

\begin{enumerate}
    \item \textbf{Image motivation:} Donors value being seen as altruistic
    \item \textbf{Payment signal:} Payment suggests donation is about money, not altruism
    \item \textbf{Net effect:} If $\beta \cdot \text{signal} > \alpha \cdot b$, payment reduces donations
\end{enumerate}

\textbf{Empirical Evidence (Mellström \& Johannesson, 2008):}
\begin{itemize}
    \item Control: 52\% donation rate
    \item Payment (\$7): 30\% donation rate (women), 53\% (men)
    \item Payment to charity: 53\% donation rate
\end{itemize}

\textbf{EBF Intervention Design:}
\begin{enumerate}
    \item \textbf{Avoid small payments:} Either no payment or substantial payment
    \item \textbf{Frame as gift/recognition:} ``Thank you gift'' vs. ``payment''
    \item \textbf{Offer charity option:} Payment to charity preserves image motivation
    \item \textbf{Emphasize social impact:} Strengthen intrinsic motivation
\end{enumerate}

\textbf{Predicted Optimal Strategy:} Recognition gifts (non-monetary) + social impact messaging + option to donate payment to charity.
\end{tcolorbox}

% ============================================================================
% SUMMARY
% ============================================================================
\section{Summary}
\label{sec:rb-summary}

\begin{tcolorbox}[colback=blue!5!white,colframe=blue!75!black,title=Key Takeaways]
\begin{enumerate}
    \item \textbf{RB-1 (Self-Confidence):} Beliefs about ability affect motivation; rational agents engage in self-deception
    \item \textbf{RB-2 (Crowding):} Incentives backfire when they signal negative information about task or trust
    \item \textbf{RB-3 (Identity):} Beliefs are assets that provide utility and enable commitment; taboos protect sacred values
    \item \textbf{RB-4 (Stratification):} Local complementarities create persistent inequality through sorting
    \item \textbf{RB-5 (POUM):} Mobility beliefs explain support for inequality even among the poor
    \item \textbf{RB-6 (Groupthink):} Organizations develop collective delusions through conformity cascades
\end{enumerate}

\textbf{Integration Principle:} Bénabou's work shows that beliefs are not just about information but serve psychological and social functions. Effective interventions must account for these functions rather than simply providing information or incentives.
\end{tcolorbox}

% ============================================================================
% GLOSSARY
% ============================================================================
\section{Glossary}
\label{sec:rb-glossary}

Key terms defined in Appendix G (REF-GLOSSARY):

\begin{itemize}
    \item \textbf{Motivated reasoning:} Biased information processing to reach desired conclusions
    \item \textbf{Self-deception:} Strategic avoidance or distortion of self-relevant information
    \item \textbf{Motivation crowding:} Extrinsic incentives reducing intrinsic motivation
    \item \textbf{POUM hypothesis:} Prospect of Upward Mobility explains tolerance for inequality
    \item \textbf{Local complementarities:} Individual outcomes depend on community composition
    \item \textbf{Groupthink:} Conformity pressure leading to collective poor decisions
    \item \textbf{Taboo:} Prohibition on considering trade-offs that would undermine sacred values
\end{itemize}

% ============================================================================
% CRITICAL FOUNDATIONS
% ============================================================================
\section{Critical Foundations}
\label{sec:rb-foundations}

\begin{enumerate}
    \item Bénabou, R., \& Tirole, J. (2002). Self-Confidence and Personal Motivation. \textit{QJE}, 117(3), 871--915.
    \item Bénabou, R., \& Tirole, J. (2003). Intrinsic and Extrinsic Motivation. \textit{RES}, 70(3), 489--520.
    \item Bénabou, R., \& Tirole, J. (2004). Willpower and Personal Rules. \textit{JPE}, 112(4), 848--886.
    \item Bénabou, R., \& Tirole, J. (2006). Incentives and Prosocial Behavior. \textit{AER}, 96(5), 1652--1678.
    \item Bénabou, R., \& Tirole, J. (2011). Identity, Morals, and Taboos. \textit{QJE}, 126(2), 805--855.
    \item Bénabou, R. (1996). Inequality and Growth. \textit{NBER Macro Annual}, 11, 11--74.
    \item Bénabou, R. (1996). Heterogeneity, Stratification, and Growth. \textit{AER}, 86(3), 584--609.
    \item Bénabou, R., \& Ok, E. (2001). Social Mobility and the Demand for Redistribution: The POUM Hypothesis. \textit{QJE}, 116(2), 447--487.
    \item Bénabou, R. (2013). Groupthink: Collective Delusions in Organizations. \textit{RES}, 80(2), 429--462.
    \item Bénabou, R., \& Tirole, J. (2016). Mindful Economics. \textit{JEP}, 30(3), 141--164.
    \item Bénabou, R., Falk, A., \& Tirole, J. (2018). Narratives, Imperatives, and Moral Reasoning. \textit{NBER Working Paper} No. 24798.
    \item \textbf{Bénabou, R., \& Henkel, L. (2025). Identity as Self-Image. \textit{Handbook of the Economics of Identity}, forthcoming. NBER WP 34297.} --- Comprehensive unified framework for belief-based identity economics.
\end{enumerate}

% ============================================================================
% OPEN ISSUES
% ============================================================================
\section{Open Issues}
\label{sec:rb-open}

\begin{enumerate}
    \item \textbf{Optimal self-deception:} How much self-deception is optimal? When does it become pathological?
    \item \textbf{Identity measurement:} How to measure identity investments and belief utility empirically?
    \item \textbf{Policy and beliefs:} How do policies shape beliefs about mobility and fairness?
    \item \textbf{Groupthink interventions:} How to design organizations that resist collective delusions?
    \item \textbf{Digital age beliefs:} How do social media and echo chambers affect motivated reasoning?
\end{enumerate}

% ============================================================================
% REFERENCES SECTION
% ============================================================================
\section{References}
\label{sec:rb-references}

\nocite{bcm_master}

This appendix integrates 33 papers from Roland Bénabou's research corpus. Full references are maintained in the master bibliography (\texttt{bibliography/bcm\_master.bib}).

\textbf{Core Bénabou Works (32 papers):}
\begin{itemize}
    \item \textbf{Motivation (with Tirole):} (2002, 2003, 2004, 2006, 2009, 2010, 2011, 2016)
    \item \textbf{Moral Reasoning (with Falk \& Tirole):} (2018) --- Narratives, imperatives, network transmission
    \item \textbf{Moral Types (with Falk \& Henkel):} (2024) --- Ends vs. Means, no stable Kantian/Utilitarian types, $\gamma \approx 0$ \textit{(NEW)}
    \item \textbf{Identity Economics (with Henkel):} (2025) --- Unified self-image framework \textit{(NEW)}
    \item \textbf{Inequality \& Growth:} (1993, 1994, 1996a, 1996b, 2000, 2002, 2005)
    \item \textbf{Political Economy:} Bénabou \& Ok (2001), Bénabou (2008)
    \item \textbf{Groupthink:} (2013, 2015)
    \item \textbf{Religion \& Innovation:} Bénabou, Ticchi \& Vindigni (2015, 2022)
    \item \textbf{Monetary Economics:} (1988, 1992), Bénabou \& Gertner (1993)
    \item \textbf{Information:} Bénabou \& Laroque (1992)
\end{itemize}

See Appendix G (REF-GLOSSARY) for term definitions and Appendix JT (LIT-TIROLE) for the Bénabou-Tirole collaboration.

\end{document}
